\documentclass[10pt,oneside]{book}
%%%% Page Info + Commands %%%%%{

%packages
\usepackage{pgfplots}
\usepackage{mathrsfs}
\pgfplotsset{compat=1.18}
\usepackage{amsmath}
\usepackage{amsfonts}
\usepackage{charter}
\usepackage{amsthm,letltxmacro}
\usepackage{amssymb}
\usepackage{fancyhdr}
\usepackage{float}
\usepackage[most]{tcolorbox}
\usepackage{enumerate}
\usepackage{enumitem}
\usepackage[dvipsnames]{xcolor}
\usepackage{changepage}
\usepackage{graphicx}

% SMALL PAGE COMMAND
\newcommand{\smallpage}{  \setlength \oddsidemargin{.25in}
            \setlength \textwidth{5.5in}
            \setlength \topmargin{0in}
            \setlength \textheight{9in}}


% Commands for sets of numbers
\newcommand{\Z}{\mathbb{Z}}\newcommand{\R}{\mathbb{R}}\newcommand{\N}{\mathbb{N}}\newcommand{\Q}{\mathbb{Q}}\newcommand{\C}{\mathbb{C}}
\newcommand{\real}[1]{\text{Re}\left(#1\right)}
\newcommand{\im}[1]{\text{Im}\left(#1\right)}

% gordie spacing and boxing commands
\newcommand{\gspc}{\vspace{0.13cm}} % 0.5 space
\newcommand{\gline}[0]{\gspc\\} % 1.5 space
\newcommand{\gbline}{\gspc\gspc\\} % double space
\newcommand{\gbox}[1]{\fbox{\parbox{\linewidth}{#1}}} % multiline box command
\newcommand{\gmod}[1]{\,(\text{mod}\,#1)}


\newcommand{\gimage}[2]{
\begin{figure}[H]
    \centering
    \includegraphics[width=#2]{#1}
\end{figure}
}

\newcommand{\centerit}[1]{{\begin{center} #1 \end{center}}}
\newcommand{\ul}[1]{\underline{#1}}

\newcommand{\prt}[2]{\frac{\partial#1}{\partial#2}}

%%%%%%%% command for graphics %%%%%%%%%%%%%

%}

% COLOR BOX

\tcbset{problembox/.style={
    colback=gray!4,
    colframe=gray!50, 
	  boxrule=0pt,
    leftrule=2.5pt, 
    left=2mm, right=2mm, top=1mm, bottom=1mm,
    boxsep=2mm,
    breakable,
    sharp corners
  }
}

\tcbset{subproblembox/.style={
    colback=gray!12,
    colframe=gray!80, 
	  boxrule=0.5pt,
    leftrule=2.5pt, 
    left=2mm, right=2mm, top=1mm, bottom=1mm,
    boxsep=2mm,
    breakable,
    sharp corners
  }
}

% PROBLEM
\newenvironment{problem}[1]
  {\vspace{-0.1cm}\begin{tcolorbox}[problembox]
   \textit{Problem. }#1}
  {\end{tcolorbox}\gspc}

  \newenvironment{subproblem}[2]
  {\vspace{-0.1cm}\begin{tcolorbox}[subproblembox]
   \textit{#1 }#2}
  {\end{tcolorbox}\gspc}


%%%% Page Setup %%%%%%%
\smallpage\sloppy
\pagestyle{fancy}
\lhead{\large{\textbf{PS2 - Maxwell's Equations}}}
\setlength{\headheight}{10pt}


% color commands
\newcommand{\cg}[1]{\color{OliveGreen}#1\color{white}}
\newcommand{\cb}[1]{\color{BlueViolet}#1\color{white}}


\newcommand{\cpr}[1]{\left(#1\right)}
%%%%%%%%%%%%% PHYSICS SPECIFIC COMMANDS

\newcommand{\vA}{\vec{\mathbf{A}}}
\newcommand{\vB}{\vec{\mathbf{B}}}
\newcommand{\vC}{\vec{\mathbf{C}}}
\newcommand{\vZ}{\vec{\mathbf{0}}}
\newcommand{\norm}{\vec{\mathbf{n}}}
\newcommand{\pvec}[1]{\left\langle#1\right\rangle}
\newcommand{\ar}{\vec{\mathbf{r}}}
\newcommand{\arhat}{\hat{\mathbf{r}}}
\newcommand{\vv}{\vec{\mathbf{v}}}
\newcommand{\delfull}{\pvec{\prt{}{x}, \prt{}{y},\prt{}{z}}}
\newcommand{\delinput}[3]{\pvec{\prt{}{x}\cpr{#1}, \prt{}{y}\cpr{#2},\prt{}{z}\cpr{#3}}}
\newcommand{\crossprod}[6]{\text{Det}\left|\begin{matrix}\hat{\textbf{i}}&\hat{\textbf{j}} &\hat{\textbf{k}}\\#1 & #2 & #3\\#4 & #5&#6\end{matrix}\right|}
\newcommand{\dps}{\displaystyle}
\newcommand{\delmatrix}[3]{\begin{bmatrix}\dps\hat x&\dps\hat y&\dps\hat z\\\dps\prt{}x&\dps\prt{}y&\dps\prt{}z\\\dps#1&\dps#2&\dps#3\end{bmatrix}}

%%%%%%%%%%%%%%%%%%%%%%%%%%%%
%%%%%%%%%%%%%%%%%%%%%%%%%%%%%%
%%%%%%%%%%%%%%%%%%%%%%%%%%%%%
%%%%%%%%%%%%%%%%%%%%%%%%%%%%%
%%%%%%%%%%%%%%%%%%%%%%%%%%%%%%
%%%%%%%%% begin doc


\begin{document}
\LetLtxMacro{\oldint}{\int}
\renewcommand{\int}[2]{{\oldint_{#1}^{#2}\!}}
\noindent Assigned Problems:\\
\textbf{Chatper 1: }29, 33, 36, 40, 43, 48, 52\\
\textbf{Chapter 2: }1
\subsection*{Chapter 1}
\subsubsection*{Problem 29}
Calculate the line integral of the function $\vv = \pvec{x^2, 2yz, y^2}$ from the origin to the point $\pvec{1, 1, 1}$ by three different routes.
\begin{enumerate}
  \item [(a)] $\pvec{0, 0, 0} \to \pvec{1, 0, 0}\to \pvec{1, 1, 0}\to \pvec{1, 1, 1}$
  \begin{problem}
    We begin by paramaterizing the curve for each part of the problem. \gline{}
    For $\pvec{0,0,0}\to\pvec{1,0,0}$, we paramaterize with $\vec r(t)= \pvec{t,0,0}$, $\vec r'(t) = \pvec{1, 0, 0}$. Solving the integral for $t:[0, 1]$ gives:
    \begin{align*}
      \int{0}   1\pvec{t^2, 0, 0}\cdot \pvec{1, 0 ,0} \,dt=\left[\cpr{\frac{t^3}3}\Big|_0^1\right]=\frac{1}3
    \end{align*}
    For $\pvec{1, 0, 0}\to \pvec{1, 1,0}$, we paramaterize as $\vec r(t) = \pvec{1, t, 0}, \vec r '(t) = \pvec{0, 1, 0}$. Solving the integral for $t:[0, 1]$ gives:
    \begin{align*}
      \int{0}   1 \pvec{1, 0, t^2}\cdot\pvec{0, 1, 0}\,dt = 0
    \end{align*}
    For $\pvec{1, 1,0}\to \pvec{1, 1, 1}$, we paramaterize as $\vec r(t) = \pvec{1, 1, t}, \vec r'(t) = \pvec{0, 0, 1}$. Solving the integral for $t:[0,1]$ gives:
    \begin{align*}
      \int{0}   1 \pvec{1, 2t, 1}\cdot \pvec{0, 0, 1}\, dt = \int{0}   1 1\,dt = 1
    \end{align*}
    Adding together all of our integrals gives:
    \begin{align*}
      \frac{1}3 + 0 + 1 = \boxed{\frac{4}3}
    \end{align*}
  \end{problem}
  \item [(b)]$\pvec{0,0,0}\to\pvec{0,0,1}\to\pvec{0, 1, 1}\to\pvec{1, 1, 1}$
  \begin{problem}
    We follow the same process as the previous question. We will paramaterize the curve for each part of the problem. \gline{}
    For $\pvec{0,0,0}\to\pvec{0,0,1}$, we paramaterize as $\vec r (t) = \pvec{0, 0, t}, \vec r'(t) = \pvec{0, 0, 1}$. For $t: 0\to 1$, we have that:
    \begin{align*}
      \int{0}   1 \pvec{0, 0, 0}\cdot\pvec{0,0, 1}\,dt= 0
    \end{align*}
    For $\pvec{0,0,1}\to\pvec{0, 1, 1}$, we paramaterize as $\vec r(t)= \pvec{0, t, 1},r'(t) = \pvec{0, 1, 0}$. For $t:0 \to 1$, we have that:
    \begin{align*}
      \int{0}   1 \pvec{0, 2t, 1}\cdot\pvec{0, 1, 0}\,dt = \int{0}   1 2t \, dt = t^2\Big|_0^1 = 1
    \end{align*}
    Finally, for $\pvec{0,1,1}\to\pvec{1,1,1}$, we paramaterize as $\vec r(t)=\pvec{t,1,1}, r'(t)=\pvec{1, 0, 0}$ for $t:0\to 1$.
    \begin{align*}
      \int{0}   1\pvec{t^2, 2, 1}\cdot\pvec{1, 0,0}\,dt=\int{0}   1 t^2\,dt = \frac{1}3t^3\Big|_0^1=\frac{1}3
    \end{align*}
    Adding up all of our integrals gives:
    \begin{align*}
      1 + 0 +\frac{1}3=\boxed{\frac{4}3}
    \end{align*}
  \end{problem}
  \item [(c)]The direct straight line.
  \begin{problem}
    We will imagine our curve as going directly from $\pvec{0, 0,0}\to\pvec{1, 1,1}$. Thus, we have our paramaterization as $\vec r(t) = \pvec{t, t, t}\, r'(t)=\pvec{1, 1, 1}$, for $t:[0, 1]$. Plugging into our integral gives:
    \begin{align*}
      \int{0}   1\pvec{t^2, 2t^2, t^2}\cdot\pvec{1, 1, 1}\,dt &= \int{0}   1 4t^2\,dt\\
      &= \frac{4}3t^3\Big|_0^1\\
      &= \boxed{\frac{4}3}
    \end{align*}
  \end{problem}
  \item [(d)] What is the line integral aournd the closed loop that goes \textit{out} along path (a) and \textit{back} along path (b)?
  \begin{problem}
    If we reverse the integrals along path $b$, we get that the line integral is just negative of our intiial result, so we get that:
    \begin{align*}
      \frac{4}3-\frac{4}3 = 0
    \end{align*}
    However, there is a more intuitive solution to this problem. Consider the potential function $f(x, y, z) = \frac{1}3x^3 + y^2z$. Note that $\nabla f(x,y,z) = \vv$. Thus, because we can generate a potential function for our vector field, our vector field is conservative.\gline{}
    Thus, the solution to this problem (and any point where the start and end paths are the same), is zero. We can check our potential function against each part of this problem:
    \begin{align*}
      f(1, 1, 1) - f(0, 0, 0) = \frac{4}3
    \end{align*}
  \end{problem}
\end{enumerate}
\newpage
\subsubsection*{Problem 33}
Test the divergence theorem for the function $\vv = \pvec{xy, 2yz,3zx}$. Take your volume as the cube shown below, with sides of length two:
\gimage{image.png}{4cm}
\begin{problem}
  To show the divergence theorem holds for this box, we will show that:
  \begin{align*}
    \iint_S \vv\cdot d\vec S \,dS = \iiint_V \nabla\cdot \vv \, dV
  \end{align*}
  First, we check the brute force way, with each side of the cube.
  \begin{subproblem}{Subproblem}
    \textbf{Bottom Side.} Consider for the bottom side of the cube, we have our surface as $S(u,v) = \pvec{u, v, 0}$, for $u, v: [0, 2]$. Now, we calcualte $S_u \times S_v$ to be:
    \begin{align*}
      \crossprod{1}{0}{0}{0}{1}{0} = \pvec{0,0, 1}
    \end{align*}
    Now, dotting our surface with our field gives:
    \begin{align*}
      \int{0}   2\int{0}   2 \pvec{uv, 0, 0}\cdot \pvec{0, 0, 1}\,du\,dv = {0}
    \end{align*}
    \textbf{Top Side.} On the top side of the cube, we model our surface as $S = \pvec{u, v, 2}$ for $u,v:[0, 2]$. We calculate our cross product to be the same as the \textbf{Bottom Side}, $\pvec{0, 0, 1}$. Plugging into our integral gives:
    \begin{align*}
      \int{0}2\int{0}2 \pvec{uv, 4v, 6u}\cdot\pvec{0, 0, 1}\,du dv &= \int 0 212\,dv\\
      &= \boxed{24}
    \end{align*}
    \textbf{Left Side.} On the left size of the box, we paramaterize our surface as $S(u,v)=\pvec{u, 0, v}$, for $u,v:[0,2]$. Calculating the cross product of $S_u, S_v$ gives us:
    \begin{align*}
      \crossprod{0}{0}{1}{1}{0}{0}=\pvec{0, -1, 0}
    \end{align*}
    Plugging our paramaterization into the surface integral gives:
    \begin{align*}
      \int 0 2\int 0 2\pvec{0, 0, 3uv}\cdot\pvec{0, 1, 0}\,dudv = 0
    \end{align*}
  \end{subproblem}
  \begin{subproblem}{}
    \textbf{Right Side}. For the side of the cube, our paramaterization is $S(u,v)=\pvec{u, 2, v}$ for $u,v:[0,2]$. The cross product is the same as above. Plugging in our paramaterization gives:
    \begin{align*}
      \int 0 2\int 0 2 \pvec{2u, 4v, 3uv}\cdot\pvec{0, 1, 0}\,dudv &= \int 0 2-8v\,dv\\
      &=-4v^2\Big|_0^2 = \boxed{16}
    \end{align*}
    \textbf{Back Side.} For the back side of the cube, we paramaterize as $S(u,v) = \pvec{0, u, v}$. Calculating the cross product gives:
    \begin{align*}
      \crossprod{0}{1}{0}{0}{0}{1}=\pvec{1, 0, 0}
    \end{align*}
    Plugging our paramaterization into the integral gives:
    \begin{align*}
      \int02\int02\pvec{0, 2uv, 0}\cdot\pvec{1, 0, 0}\,dudv = \boxed{0}
    \end{align*}
    \textbf{Front Side.} For the front side of the cube, we paramaterize with $S(u,v)=\pvec{2, u, v}$. The cross product is the same as above. Plugging into the surface integral gives:
    \begin{align*}
      \int02\int02\pvec{2u, 2uv, 6v}\cdot\pvec{1, 0, 0}\,dudv &= \int02 4 \,dv=\boxed{8}
    \end{align*}
    Adding up all of our results together gives:
    \begin{align*}
      0 + 0 + 0 + 24 + 8 + 16 = \boxed{48}
    \end{align*}
  \end{subproblem}
  Now, we check if this aligns with our expectations for the divergence theorem. 
  \begin{subproblem}{Subproblem.}
    We want to check to see if $\iiint_V \nabla\cdot \vv\,dV = 16$. Our cube follows the simple paramaterization: $V(x,y,z)=\pvec{x,y,z}$, for $x,y,z:[0, 2]$. Now, we calculate the divergence:
    \begin{align*}
      \nabla\cdot \vv = y+2z+3x.
    \end{align*}
    Plugging into our equation gives us:
    \begin{align*}
      \int02\int02\int02 3x + y + 2z\,dxdydz &= \int02\int02 6 + 2y + 4z\,dydz\\
      &= \int02 12 + 4 + 8z\,dz\\
      &= 24 + 8 + 16 = \boxed{48}
    \end{align*}
  \end{subproblem}
  Thus, the divergence theorem holds in this case, and is significantly easier than taking a total of six surface integrals. 
\end{problem}
\subsubsection*{Problem 36}
\begin{enumerate}
  \item [(a)] Show that:
  \begin{align*}
    \int S{}f(\nabla\times \vA)\cdot d\vec a = \int S{}[\vA \times(\nabla f)]\cdot d\vec a + \oint_P f\vA \cdot d\vec\ell
  \end{align*}
  \begin{problem}
    Consider that via the product rule, we have that:
    $$\nabla \times (f\vA)=f(\nabla\times A )- \vA \times(\nabla f).$$
    Thus, we have that:
    \begin{align*}
      \int S{}f(\nabla\times\vA)\cdot d\vec a = \int S{}\nabla \times (f\vA) +  \vA \times(\nabla f)\cdot d\vec a
    \end{align*} 
    Note that via Stoke's theorem, we have that:
    $$\int S{} \nabla\times(f\vA)=\oint_P f\vA \cdot d\vec\ell$$
    Thus, we have the relation:
    \begin{align*}
      \int S{}f(\nabla\times \vA)\cdot d\vec a = \int S{}[\vA \times(\nabla f)]\cdot d\vec a + \oint_P f\vA \cdot d\vec\ell.
    \end{align*}
  \end{problem}
  \item [(b)] Show that:
  \begin{align*}
    \int V{} \vB\cdot(\nabla\times\vA)d\tau = \int V{}\vA\cdot(\nabla\times\vB)d\tau + \oint_S(\vA\times\vB)\cdot d\vec a
  \end{align*}
  \begin{problem}
    Via one of the product rules, we have that: 
    $$\nabla\times(\vA\times\vB)=\vB\cdot(\nabla\times \vA)-\vA\cdot(\nabla F).$$
    Thus, we have that:
    \begin{align*}
      \int V{}\vB\cdot(\nabla\times\vA)d\tau = \int V{}\nabla\cdot(\vA\times\vB)+\vA\cdot(\nabla\times \vB)
    \end{align*}
    I notice something resembling the divergence theorem, so we substitute in for our answer:
    \begin{align*}
      \int V{} \vB\cdot(\nabla\times\vA)d\tau =  \oint_S(\vA\times\vB)\cdot d\vec a+\int V{}\vA\cdot(\nabla\times\vB)d\tau 
    \end{align*}
  \end{problem}
\end{enumerate}
\subsubsection*{Probelm 40}
Compute the divergence of the function:
\begin{align*}
  \vv = \pvec{r\cos\theta, r\sin\theta, r\sin\theta\cos\phi}
\end{align*}
Check the divergence theorem for this function, using as your volume the inverted hemispherical bowl of radius $R$, resting on the $xy$ plane and centered at the origin. 
\begin{problem}
  We will compute the divergence, and check the divergence theorem for this function. 
  \begin{subproblem}{Solving the Divergence.}
    First, we check the divergence. Note that our $\nabla$ is different for spherical:
    \begin{align*}
      \nabla = \pvec{\prt{}r, \quad\frac{1}{r}\prt{}{\theta},\quad \frac{1}{r\sin\theta}\prt{}{\phi}}
    \end{align*}
    Now, we begin solving:
    \begin{align*}
      \nabla\cdot\vv &= \frac{1}{r^2}\prt{}r(r\,r\cos\theta)+\frac{1}{r\sin\theta}\prt{}{\theta}(r\,r\sin\theta\sin\theta) + \frac{1}{r\sin\theta}\prt{}{\phi}(r\sin\theta\cos\phi)\\
      &=5\cos\theta - \sin\phi
    \end{align*}
  \end{subproblem}
  Next, we solve the divergence theorem here:
  \begin{subproblem}{Checking Divergence Theorem. }
    So first we'll prove the $\int V{}$ part of the problem:
    \begin{align*}
      \int V{}(\nabla\cdot\vv)d\tau &= \int0{\pi/2}\int0{2\pi}\int0{R}(5\cos\theta - \sin\phi)r^2\sin\theta \,drd\theta d\phi\\
      &= \int0{\pi/2}\int0{2\pi}\frac{1}3R^3(5\cos\theta - \sin\phi)r^2\sin\theta\\
      &\to \text{sin goes to 0}\\
      &\Rightarrow \cos\theta\sin\theta = \frac{1}2\sin(2\theta)\\
      &= -\frac{10\pi}{3}R^3\frac{1}4\cos\theta \Big|_0^{\pi/2}\\
      &= \frac{5\pi}{3}R^3
    \end{align*}
    Next, we'll show that this works with the surface integral. 
  \end{subproblem}
\end{problem}
\subsection*{Problem 43}
\begin{enumerate}
  \item Find the divergence of the function:
  \begin{align*}
    \vv = \pvec{s(2+\sin^2\theta), s\,\sin\phi\cos\phi, 3z}
  \end{align*}
  \begin{problem}
    Note that our divergence formula is a little bit different:
    \begin{align*}
      \nabla\cdot f =\pvec{ \frac{1}{s}\prt{f_s}{s},\frac{1}s\prt{f_\phi}\phi,\prt{f_z}{z} }
    \end{align*}
    Now, we solve for the divergence. 
    \begin{align*}
      \nabla\cdot \vv &= \frac{1}s\cdot\prt{s(2+\sin^2\theta)}{s} + \frac{1}s\prt{s^2\sin\phi\cos\phi}{\phi}+\prt{3z}{z}\\
      &= 2(2+\sin^2\theta)+(\cos^2\theta - \sin^2\theta)+3\\
      &= 4 + 2\sin^2\theta - \sin^\theta +\cos^2\theta + 3\\
      &= 8
    \end{align*}
  \end{problem}
  \item Now, we have to test for the divergence theorem. 
  \begin{problem}
    First, we apply the actual divergence theorem.
    \begin{subproblem}{Divergence Theorem.}
      Here, because we just have a constant, 8, for our divergence, the integral simplifies to:
      \begin{align*}
        8\cdot\iiint_V 1 dV
      \end{align*}
      So this is just the area of the quarter cylinder, which is $\frac{1}4(5*4\pi)\cdot 8 = 8*5\pi = 40\pi$
    \end{subproblem}
    Next, we test for each surface fo the cylinder. 
    \begin{subproblem}{Solving Surface Integrals.}\gline{}
      \textbf{Bottom: } Paramaterization: $\pvec{2, t, 0}, '\to\pvec{0,1,0}$ $t:[0,\pi/2]\Rightarrow \vv\cdot d\vec a = 0$.
      \textbf{Top: }$\pvec{2, t, 5} '\to\pvec{0, 1,0}, t:[0, \pi/2]\Rightarrow 15\cdot\int02 s\int0{\pi/2}d\phi=15\pi$ \\
      \textbf{Front:} Mathematics gives $25\pi$.\\
      \textbf{Left and back side: }0\gline{}
      Total: $40\pi$.
    \end{subproblem}
    Thus, the divergence theorem holds
  \end{problem}
  \item [(c)] The curl here will be 0. 
\end{enumerate}
\subsection*{Problem 48}
Evaluate the following integrals:
\begin{enumerate}
  \item [(a)] $\int{}{} (r^2+\vec r \cdot \vec a + a^2)\delta^3(\vec r - \vec a)$
  \begin{problem}
    So we evaluate the value of $\vec r$ at $\vec a$. Thus, we just get:
    \begin{align*}
      \int{}{}a^2 + \vec a \cdot \vec a + a^2 \,dr = 3a^2
    \end{align*}
  \end{problem}
  \pagebreak
  \item [(b)] $\int{}{}(\vec r - \vec b)^2\frac{1}{125}\delta^3(5\vec r)d\tau$
  \begin{problem}
    All we do is evaluate it when $\vec r=\vec0$.This just gives us:
    \begin{align*}
      (\vec b^2)\frac{1}{125}.
    \end{align*}
    Because $\vec b = \pvec{4, 3}$, we have that $|b|^2 =25$, so we just get a result of $\boxed{\frac{25}{125}=\frac{1}5}$.
  \end{problem}
  \item [(c)]$\int{V}{}[r^4+r^2(\vec r \cdot \vec c)+c^4]\delta^3(\vec r - \vec c)d\tau$, sphere of radius 6, for $\vec c = \pvec{5, 3, 2}$.
  \begin{problem}
    Thus, we evaluate $\vec r$ at $\vec c$. We get:
    \begin{align*}
      \int{V}{}{c^4 + c^2(\vec c \cdot \vec c) + c^4}\,d\tau = \int{V}{}3c^4\,d\tau.
    \end{align*}
    Note that $c^4$ here is $1444$. This is actually outside of our sphere, so the resulting integral is probably zero. I'm not gonna check.  
  \end{problem}
  \item [(d)]$\int{V}{}\vec r \cdot (\vec d - \vec r)\delta^3(\vec e - \vec r)d\tau$. Here, $\vec d = \pvec{1,2, 3}, \vec e = \pvec{3, 2, 1}$, $V$ is a sphere of radius 1.5 at $\pvec{2, 2, 2}$. 
  \begin{problem}
    We begin by evaluating the integral at $\vec e = \vec r$. Thus, we have:
    \begin{align*}
      \vec e \cdot(\vec d - \vec e)
    \end{align*} 
    Thus, we just evaluate and get:
    \begin{align*}
      \pvec{3, 2, 1}\cdot\pvec{-2,0,2}= -4
    \end{align*}
  \end{problem}
\end{enumerate}
\subsubsection*{Problem 52}
For Theorem 2, show that $(d)\Rightarrow (a), (a)\Rightarrow (c), (c)\Rightarrow(b), (b)\Rightarrow (c), $and $(c)\Rightarrow (a)$.  
\begin{problem}
  \begin{subproblem}
    First we show $(d)\Rightarrow(a)$.\gline{}
    Consider that $F$ is the curl of a vector function. $F=-\nabla \vA$. Thus, we have that:
    \begin{align*}
      \nabla\cdot(\nabla\times \vA)
    \end{align*}
    Note that with our $\nabla\cdot \vA$, we have that:
    \begin{align*}
      \nabla\times\vA=\pvec{\prt{a_2}{z}-\prt{a_3}{y}, \prt{a_1}{z} - \prt{a_3}{a}, \prt{a_2}{x}-\prt{a_1}{y}}
    \end{align*}
    Then, because partials distribute with $\nabla$, we basically just end up with a bunch of terms that cancel, like for:
  \end{subproblem}
  \begin{subproblem}{}
    \begin{align*}
      \nabla = \pvec{\prt{}{x},\prt{}{y},\prt{}{z}}
    \end{align*}
    We get $\nabla\cdot(\nabla\times\vA)=\vec 0$
  \end{subproblem}
  \begin{subproblem}{}
    Next, we show $(a)\Rightarrow(c)$.\gline{}
    We assume $\nabla\cdot F = 0$.This one's pretty simple. Via the divergence theorem, we convert:
    \begin{align*}
      \oint F \cdot d\vec a = \iiint_V \nabla\cdot F = 0
    \end{align*}
  \end{subproblem}
  \begin{subproblem}{}
    Next,we show that $(c)\Rightarrow(b)$. \gline{}
    We assume that $\oint F\cdot d\vec a=0$. Note that for any given surface, we know that via green's theorem, that if we combine any number of line integrals that go back to the same point, we get:
    \begin{align*}
      \int{I_1}{}F\cdot d\vec a + \int{I_2}{}F\cdot d\vec a \cdot  + \cdots =\oint F\cdot d\vec a = 0
    \end{align*}
    Thus, regardless of surface, boundary line is identical. 
  \end{subproblem}
  \begin{subproblem}{}
    Now, we show that $(c)\Rightarrow (b)$. This is the same exact proof as the last problem, beacuse it's an equality proof, so it's automatically an ``if an only if'' proof. 
  \end{subproblem}
  \begin{subproblem}{}
    Next, we will show $(a)\Rightarrow(c)$. The proof from $(c)\Rightarrow(a)$ is also an equality proof, so $(c)\Rightarrow (a) \implies (a)\Rightarrow(c)$ for my example.  
  \end{subproblem}
\end{problem}
\subsection*{Chapter 2}
\subsubsection*{Problem 1}
\begin{enumerate}
  \item [(a)] Twelve equal charges, $q$, are situated at the corners of a regular 12-sided polygon. What is the net force on a test charge $Q$ at the center. 
  \begin{problem}
    The net force on this test charge is 0, because for each charge, there is another equal charge equidistance, but opposite in unit vector from it, so thus, force cancels. 
  \end{problem}
  \pagebreak
  \item[(b)] Suppose \textit{one} of the 12 $q$s is removed (at 6 o-clock). What is the force on $Q$? Explain your reasoning.
  \begin{problem}
    The same cancelation applies to all the charges in the previous problem, \textit{except} the one at 12 o'clock, so the unit vector will point directly down from it:
    \begin{align*}
      \vec F = \frac{1}{4\pi\varepsilon_0}\cdot\frac{Qq}{r^2}\cdot \hat r
    \end{align*}
    Here, $\hat r = \pvec{0, -1, 0}$ if we assume our polygon is bound in the x-y plane, and $r$ is the distance from center to one point of the polygon. 
  \end{problem}
  \item[(c)] Now 13 equal charges, $q$, are placed at the corners of a regular 13-sided polygon. What is the force on a test charge $Q$ at the center?
  \begin{problem}
    Now, this one appears a little more complicated, but I think it's still just $\textbf{zero}$. Every charge pulls equally on the charge at fixed intervals of angles around $Q$. Thus, if the net force was \textit{not} zero, then there is a charge that the Force \textit{leans} towards. However, this is not possible because all charges are indistinguishable from one another. Thus the net charge must be zero. 
  \end{problem}
  \item[(d)] If one of the 13$q$s is removed, what is the force on $Q$? 
  \begin{problem}
    Instead of treating these as charge pairs, imagine we just placed a charge of $-q$ at the place where we removed that $q$ instead of removing it. It's an equivalent system in terms of electric field. Thus, we just point towards that charge:
    \begin{align*}
      \vec F = \frac{1}{4\pi\varepsilon_0}\cdot\frac{Qq}{r^2}\cdot\hat r
    \end{align*}
    Where $\hat r$ points towards the 'missing' charge. 
  \end{problem}
\end{enumerate}
\end{document}